\section{Analytic and independent reference benchmarks}
\label{sec:analytic-benchmarks}

This section freezes the minimum benchmark equations required before neural
training. Exact formulas are used only within their known asymptotic/domain
assumptions.

\subsection{Uniform DC conductor}

For a straight homogeneous conductor of length $L$, conductivity $\sigma$, and
cross-sectional area $A$,
\[
  R_{\rm dc}
  =
  \frac{L}{\sigma A}.
  \label{eq:dc-resistance}
\]
For a superelliptic cross-section
\[
 \left|\frac{x}{a}\right|^m+
 \left|\frac{y}{b}\right|^m\le1,
\]
the area is
\[
 A
 =
 4ab\,
 \frac{\Gamma(1+1/m)^2}
      {\Gamma(1+2/m)}.
 \label{eq:superellipse-area}
\]
The low-frequency conductor solver must converge to
Equation~\eqref{eq:dc-resistance} using the same geometry and port convention.

\subsection{Circular-loop mutual inductance}

For two coaxial filamentary circular loops with radii $a,b$ and axial
separation $z$, define
\[
  k^2
  =
  \frac{4ab}{(a+b)^2+z^2}.
\]
The magnetoquasistatic mutual inductance is
\[
  M
  =
  \mu\sqrt{ab}\,
  \frac{(2-k^2)K(k)-2E(k)}{k},
  \label{eq:coaxial-mutual}
\]
where $K,E$ are complete elliptic integrals. In the thin-conductor,
electrically-small limit, the off-diagonal reactance must satisfy
\[
  \Ima Z_{12}
  \approx
  \omega M
\]
under the fixed $e^{+j\omega t}$ convention.

This benchmark is evaluated away from coincident-loop singular limits.

\subsection{General filament Neumann integral}

For two disjoint thin closed curves $C_1,C_2$ in the MQS limit,
\[
  M_{12}
  =
  \frac{\mu}{4\pi}
  \oint_{C_1}\oint_{C_2}
  \frac{\dd\mathbf l_1\cdot\dd\mathbf l_2}
       {\norm{\mathbf r_1-\mathbf r_2}}.
  \label{eq:neumann-mutual}
\]
High-order independent curve quadrature of
Equation~\eqref{eq:neumann-mutual} provides a pose-general benchmark for the
leading mutual coupling of sufficiently thin conductors. It is independent of
the finite-cross-section IE implementation.

\subsection{High-frequency round-wire resistance asymptote}

For a long circular conductor of radius $a$ with
\[
  a/\delta\gg1,
\]
current is concentrated in a surface layer. The leading AC resistance per unit
length is
\[
  R'_{\rm ac}
  \sim
  \frac{1}{2\pi a\sigma\delta}.
  \label{eq:skin-asymptote}
\]
The finite-cross-section modal solver must approach this trend as conductor
curvature/proximity effects are removed and $a/\delta$ increases.

This is an asymptotic benchmark, not an exact finite-frequency formula.

\subsection{Rigid-motion invariance}

For any common $(R_0,t_0)\in\SE(3)$ in an isotropic homogeneous background,
\[
  Z(g_0\cdot\mathcal S)
  =
  Z(\mathcal S).
  \label{eq:rigid-benchmark}
\]
The test suite applies deterministic random rotations and translations over
many orders of orientation representation, comparing both port observables and
equivariantly transformed continuous field/loss queries.

\subsection{Reciprocity benchmark}

For reciprocal materials and consistently oriented ports,
\[
  Z-Z^\T=0.
\]
Report
\[
  \eta_{\rm rec}
  =
  \frac{\norm{Z-Z^\T}_F}
       {\max(\norm Z_F,Z_0)}.
\]

\subsection{Power benchmark}

For each deterministic complex current vector $i$,
\[
 P_{\rm port}
 =
 \frac12 i\adj\Herm(Z)i.
\]
The independently integrated conductor loss is
\[
 P_c
 =
 \frac12
 \int_{\Omega_c}
 \mathbf J\adj\rho_e\mathbf J\,\dd V.
\]
For an electrically-small non-radiating benchmark,
\[
  P_{\rm port}\approx P_c.
\]
For a retarded/radiating benchmark,
\[
  P_{\rm port}=P_c+P_{\rm out}
\]
with $P_{\rm out}$ computed from the independent radiation channel.

\subsection{Infinite-medium thermal impulse}

For homogeneous infinite material with diffusivity
\[
  \alpha=\frac{k}{\rho c},
\]
an instantaneous point heat input of energy $Q$ at the origin gives, for $t>0$,
\[
 \Delta T(r,t)
 =
 \frac{Q}
      {\rho c(4\pi\alpha t)^{3/2}}
 \exp\left(
  -\frac{r^2}{4\alpha t}
 \right).
 \label{eq:thermal-impulse}
\]
This tests transient Green-kernel normalization and time scaling.

\subsection{Infinite-medium thermal step source}

For a point source of constant power $P$ switched on at $t=0$,
\[
  \Delta T(r,t)
  =
  \frac{P}{4\pi kr}
  \operatorname{erfc}
  \left(
    \frac{r}{2\sqrt{\alpha t}}
  \right),
  \qquad r>0.
  \label{eq:thermal-step}
\]
As $t\to\infty$,
\[
  \Delta T(r,\infty)=\frac{P}{4\pi kr}.
\]
This benchmark covers short-time, transition, and steady limits of the thermal
transfer.

\subsection{Lumped thermal mode}

For a single thermal capacity $C$ connected to ambient by conductance $G$,
\[
  C\dot\theta+G\theta=P,
\]
with $\theta(0)=0$ and constant $P$,
\[
 \theta(t)
 =
 \frac{P}{G}
 \left(
  1-e^{-Gt/C}
 \right).
 \label{eq:lumped-thermal}
\]
The reduced integrator and steady solver must reproduce this exactly to their
declared time-integration tolerance.

\subsection{Independent numerical oracle}

Analytic formulas do not cover finite-cross-section arbitrary superellipse
coils in generic poses. A locked subset is therefore compared with an
independent converged numerical method. The legacy volume H(curl) solver may be
used only after its own mesh/domain/self-physics convergence is demonstrated for
that benchmark. It is not considered independent evidence when configured from
the same integral basis or quadrature code.

\subsection{Benchmark tiers}

Each benchmark is labeled:
\begin{description}
\item[A] analytic/exact under explicitly stated assumptions;
\item[B] asymptotic with a declared small parameter;
\item[C] independent numerical cross-method;
\item[D] structural identity such as reciprocity, rigid invariance, or energy
  closure.
\end{description}

REFERENCE acceptance requires all applicable Tier A/D tests, selected Tier B
regime tests, and at least one Tier C cross-method family before neural labels
are trusted.
